\section{The Hyperbolic Bulk Toolkit}

The mesh is a $\{3,4,3,4\}$ honeycomb, and that honeycomb lives in four-dimensional hyperbolic space. Measuring physics inside such a space is hard. The space curves away from itself exponentially, so the naive habits that work on a flat lattice quietly break. A sum over a finite chunk of the mesh, the kind that reads off a Green's function or a spectral density on an ordinary crystal, comes back contaminated rather than wrong by a little. This appendix is the field guide to the tools that get clean answers from the curved bulk. For each tool it names the exact problem the tool dissolves and the exact number the tool delivers.

Two regions of the mesh behave very differently, and the toolkit respects that split. The deep interior, the \emph{bulk}, is genuinely hyperbolic, with its exponential room. The thin flat collar at infinity, the \emph{cusp}, is genuinely flat and three-dimensional. Different geometry needs different instruments. The bulk gets the recursion, the closed manifold, the band picture, and the tensor network. The cusp gets the difference method. Both rest on two careful builders that lay down the docks in the first place.

\subsection{Why the naive methods fail}

A finite piece of hyperbolic space is almost all skin. In flat three-space a ball of radius $r$ holds a volume that grows like $r^3$ while its surface grows like $r^2$, so the interior eventually swamps the boundary. In hyperbolic space the reverse holds and holds violently. Volume grows exponentially, like $e^{(D-1)r/R}$ with $R$ the curvature radius, and so does the bounding surface. The surface never falls behind the interior. A finite patch of the bulk is essentially boundary all the way through.

\begin{principle}[The all-boundary problem]
In hyperbolic space the surface of any finite patch is the same exponential order as its interior. There is no regime in which the interior dominates. A finite-patch measurement of a bulk quantity is therefore dominated by surface effects rather than by the bulk it was meant to probe.
\end{principle}

A second contamination rides along with the first. Any finite patch carries a lowest spectral mode, the constant or $k=0$ zero mode, and on an exponentially open space that mode does not decouple. It leaks into Green's functions and spectral sums and sets a false floor. A naive bulk gravity measurement that ignores both effects reads the wrong answer with confidence. An early measurement of the bulk gravity exponent on a finite patch returned roughly $-4.3$, pure garbage, produced by exactly two mistakes layered together. The all-boundary problem, and a confusion between straight-line distance through the embedding picture and true geodesic distance along the mesh. Clean bulk physics needs methods that either cancel these two contaminations outright or sidestep the finite patch altogether.

\subsection{The four bulk tools}

\subsubsection{Bethe-lattice cavity recursion}

The hyperbolic tree is self-similar. Cut one branch off and what remains is a smaller copy of the same tree. That exact self-similarity turns the resolvent, the object whose imaginary part is the spectral density and whose poles are the correlator, into a fixed point of a short recursion. This is the cavity, or self-energy, method. It removes one site, writes the self-energy felt at that site in terms of the self-energies of its neighbours on the trees that remain, and closes the loop on itself.

\begin{definition}[Cavity recursion]
The \emph{cavity recursion} solves the resolvent of the infinite Bethe lattice by a fixed-point equation for the self-energy. Each branch contributes a copy of the same self-energy, so the infinite tree is summed exactly through one self-consistent unknown, with no finite patch ever constructed.
\end{definition}

The strength of the method is that it never builds a patch at all. It works on the infinite tree directly, and it measures separation by tree distance, the true count of docks along the path, never by an embedded straight line. Both naive contaminations are absent by construction. There is no boundary to swamp the interior, and there is no spurious distance to confuse the decay. The recursion returns the bulk correlator exactly, and that correlator decays as $1/r^2$. This is the holographic decay law of the four-dimensional bulk, and the recursion is the gold-standard tool that fixes it. It is validated against a directly solved finite tree, where the two methods agree.

\begin{result}[The bulk correlator]
The cavity recursion gives the bulk two-point function exactly, with a decay
\[
G(r) \;\sim\; \frac{1}{r^2}.
\]
This is the holographic correlator of the four-dimensional hyperbolic bulk. It is read with no finite patch and with true tree distance, so neither the all-boundary nor the zero-mode contamination can touch it.
\end{result}

\subsubsection{Closed hyperbolic manifolds}

If the trouble is the boundary, the cleanest fix is to abolish the boundary. Quotient the infinite hyperbolic space by a discrete group of symmetries and the result is a finite manifold with no surface at all. It wraps onto itself the way a flat torus wraps a flat plane, except the local geometry stays hyperbolic everywhere. With no surface there is no surface contamination, and the spectra come out clean.

\begin{definition}[Closed hyperbolic manifold]
A \emph{closed hyperbolic manifold} is the quotient of hyperbolic space by a discrete fixed-point-free group. It is finite in extent and has no boundary, while its curvature stays constant and negative throughout.
\end{definition}

The worked example is a small Cayley graph on $168$ docks, of degree three, with no boundary anywhere. The $168$ is the order of the simple group $\mathrm{PSL}(2,7)$, and the symmetry that closes the space also organizes the spectrum. The Laplacian block-diagonalizes along the irreducible representations of the group, so its eigenvalues fall into clean bands whose multiplicities are exactly the irrep dimensions. This boundaryless object cross-checks two readings at once. It confirms the spectral dimension of the bulk, and it exhibits the irrep-dimension band structure that a finite patch could never show, because a finite patch is drowned by its roughly ninety-nine-percent boundary before any band can be seen.

\subsubsection{Band theory on the bulk}

Where the mesh carries translation-like symmetries, the whole infinite computation collapses onto one small unit cell. This is the Bloch picture, the same trick that turns a crystal of countless atoms into a band diagram read off a single cell. Dispersion relations and spectral gaps follow from diagonalizing a small matrix indexed by the cell, parametrized by a crystal momentum, rather than from any giant sum over docks.

\begin{principle}[Bands from a small cell]
A translation-like symmetry of the mesh lets the spectrum be read off one unit cell, parametrized by a momentum, instead of off a large patch. The dispersion and the gaps emerge from a small matrix, trading an intractable computation for a tractable one.
\end{principle}

This is the cheapest of the bulk tools and the one that scales. It does not so much defeat the all-boundary problem as refuse to meet it, since there is no patch to have a boundary.

\subsubsection{Tensor networks, the holographic contraction}

The last bulk tool builds the holographic state directly and contracts it layer by layer, from the deep interior outward toward the cusp. Each radial shell is one layer of the network, and contracting one more shell is one step of coarse-graining. This is the explicit form of the slogan that radial motion in the bulk is renormalization on the boundary. The contraction yields the same bulk correlator the recursion found, with a measured decay of $1/r^{1.93}$, matching the recursion's $1/r^2$ to within the fit.

\begin{result}[Two methods, one law]
The tensor-network contraction returns a bulk correlator decaying as $1/r^{1.93}$, agreeing with the cavity recursion's exact $1/r^2$. Two independent algorithms, one a self-consistent recursion on the infinite tree and the other an explicit layer-by-layer contraction, landing on the same decay is the strongest support for the bulk gravity law.
\end{result}

The agreement matters because the two methods share almost nothing in their machinery. The recursion never builds the state and works by self-consistency. The network builds the state explicitly and works by contraction. When they meet on the same exponent, the exponent is real and not an artifact of either method.

\subsection{The difference method for the cusp}

The cusp is flat and three-dimensional, so its Green's function should be the clean Newtonian $1/r$ of flat three-space. The naive finite-patch cusp sum, however, is wrecked by the same $k=0$ zero-mode contamination that haunts the bulk. The constant mode sets a false floor, and the measured decay sits above the true $1/r$.

The fix is to subtract the contamination away. Rather than measuring $G(r)$, the difference method measures $G(r) - G(r_0)$ against a fixed reference radius $r_0$. The constant mode is common to both terms and cancels in the difference. What remains is a regular cosine integrand with no zero-mode pathology, and it integrates to the clean flat-space result.

\begin{definition}[The difference method]
The \emph{difference method} measures the regulated quantity $G(r) - G(r_0)$ rather than $G(r)$ alone. The contaminating $k=0$ mode is identical in both terms and cancels in the subtraction, leaving a regular integrand that reports the true Green's function.
\end{definition}

\begin{result}[The flat cusp law]
After subtraction the cusp Green's function decays as $1/r$ with coefficient near $1/(4\pi)$, the exact Newtonian form of flat three-space. This supersedes the earlier confounded cusp measurements, which the uncancelled zero mode had pushed off the true value.
\end{result}

\subsection{What the toolkit fixed}

These tools are not decoration. Each one overturned a specific earlier artifact, and naming the artifacts is the clearest way to see what the tools are for.

The naive finite-patch bulk gravity exponent read roughly $-4.3$, garbage produced by the straight-line versus geodesic confusion stacked on the all-boundary problem. The cavity recursion, which builds no patch and uses true tree distance, gives the clean $1/r^2$. The tensor-network contraction, sharing none of the recursion's machinery, independently agrees.

The finite patch of the bulk is roughly ninety-nine percent boundary, so any spectral structure is drowned before it can appear. The closed $168$-dock manifold removes the boundary at the root, and on it the spectrum splits into clean bands whose multiplicities match the irrep dimensions exactly.

The cusp dimension reading does not rest on a single estimator. The spectral dimension says the cusp is three-dimensional, and a second, independent estimator drawn from the density of causal relations among the docks corroborates it. Two independent estimators agreeing, rather than one standing alone, is what licenses the claim that the cusp is genuinely flat three-space.

\subsection{The builders}

Before any of these tools can run, the docks have to be laid down, and the two regions need two different builders. Getting the builder wrong is its own source of garbage, so the split is deliberate.

\begin{principle}[Two regions, two builders]
The hyperbolic bulk and the flat cusp are built by different procedures. The hyperbolic builder projects from a disk model and works only in negative curvature. The Euclidean builder lays down an integer lattice and serves the flat controls and the flat cusp.
\end{principle}

The hyperbolic cell-direct builder grows the mesh by a breadth-first walk, placing each new dock by a projection from the disk model of hyperbolic space. It is hyperbolic only. Pointed at a flat honeycomb it stalls, because the flat cell centers project to ideal points at infinity and the walk has nowhere to step. This is a feature, not a bug. The builder refuses to pretend a flat space is curved.

The Euclidean-lattice builder is the flat counterpart, used for the flat controls and the flat cusp. It lays the integer lattice down directly, the four-dimensional $D_4$ or square lattice, or the three-dimensional cubic one, choosing by the symbol it is given. Where the hyperbolic builder would stall, this builder simply places the docks on their lattice and proceeds.

\subsection{Summary}

\begin{table}[ht]
\centering
\begin{tabular}{@{}p{3.3cm}p{4.6cm}p{4.2cm}@{}}
\toprule
tool & the problem it solves & what it yields \\
\midrule
Bethe cavity recursion & all-boundary contamination, builds no patch & exact bulk correlator, $1/r^2$ \\
closed hyperbolic manifold & surface contamination & boundaryless spectra, irrep-dimension bands \\
band theory & the cost of a huge patch & dispersion and gaps from a small cell \\
tensor-network contraction & confirming the bulk law & $1/r^{1.93}$, the radial-renormalization structure \\
the difference method & the $k=0$ cusp contamination & the clean flat-cusp $1/r$ \\
Euclidean-lattice builder & the hyperbolic builder stalls on flat & the flat controls and the flat cusp \\
\bottomrule
\end{tabular}
\caption{The bulk toolkit, the contamination each tool removes, and the number it delivers.}
\label{tab:bulk-toolkit}
\end{table}

These tools deliver the numbers used where the paper measures gravity, the cusp $1/r$ and the bulk $1/r^2$, and where it measures holography, the area law and the bulk correlator. The computational engines they sit on are described in the appendix on the computational toolkit.
